\section{Architecture}
\label{sec:architecture}

\subsection{Guiding principles}

Four principles shaped the design. They are stated here because most of the decisions described
in the following sections derive from them.

\begin{description}[leftmargin=1.2em,style=nextline]
    \item[A pure model.] The classes that \emph{are} the automaton (grid, neighborhoods, rules,
          automaton) know nothing about files, threads or pixels. They form a small library that can be
          tested in isolation and reused behind any front end.
    \item[Dependencies point one way.] Packages form layers (\cref{fig:layers}); a package only
          depends on the layers below it. Swing is confined to a single package: deleting it leaves a
          working batch application.
    \item[Immutability by default.] Positions, shapes, neighborhoods, syntax trees, definitions and
          snapshots are immutable. Mutable state is concentrated in three classes: \code{GridCell},
          \code{CellularAutomaton} and \code{Simulation}.
    \item[Errors are for the user.] Whatever comes from outside is validated at the boundary, and a
          rejected input comes with a message that says what to fix. Inner code then trusts its inputs.
\end{description}

\begin{figure}[H]
    \centering
    \begin{tikzpicture}[
            box/.style={draw=ink!60, fill=classfill, rounded corners=2pt, minimum height=0.85cm, text height=1.6ex, text depth=0.5ex,
                    font=\ttfamily\small, inner xsep=7pt},
            layer/.style={font=\sffamily\scriptsize\itshape, text=muted, anchor=east},
            dep/.style={-{Straight Barb[length=3pt]}, draw=muted, line width=0.5pt}]
        \node[box] (cli) at (2.0, 4.5) {cli};
        \node[box] (swing) at (5.0, 4.5) {ui.swing};
        \node[box] (text) at (8.2, 4.5) {ui.text};
        \node[box] (sim) at (2.6, 3.0) {simulation};
        \node[box] (render) at (6.0, 3.0) {render};
        \node[box] (explore) at (9.3, 3.0) {explore};
        \node[box, minimum width=9.6cm] (config) at (5.6, 1.5) {config};
        \node[box] (automaton) at (1.6, 0) {automaton};
        \node[box] (rule) at (4.4, 0) {rule};
        \node[box] (nb) at (7.4, 0) {neighborhood};
        \node[box] (grid) at (10.2, 0) {grid};
        \draw[dep] (automaton) -- (rule);
        \draw[dep] (rule) -- (nb);
        \draw[dep] (nb) -- (grid);
        \node[layer] at (0.2, 4.5) {entry points};
        \node[layer] at (0.2, 3.0) {services};
        \node[layer] at (0.2, 1.5) {input files};
        \node[layer] at (0.2, 0) {model};
        \begin{scope}[on background layer]
            \foreach \y in {4.5, 3.0, 1.5, 0}
            \fill[packfill] (0.35, \y - 0.62) rectangle (11.5, \y + 0.62);
            \foreach \y in {3.75, 2.25, 0.75}
            \draw[dep] (5.9, \y + 0.13) -- (5.9, \y - 0.13);
        \end{scope}
    \end{tikzpicture}
    \caption{The packages, in layers. A package depends only on packages of the layers below it
        and, within the model, on the packages to its right.}
    \label{fig:layers}
\end{figure}

\subsection{Overview of the classes}

\Cref{fig:overview} is the class diagram of the project, reduced to the classes that carry the
design and to the relationships between them. Members are omitted here; the diagrams of
\cref{sec:grid,sec:neighborhoods,sec:rules,sec:simulation} give them package by package.

\begin{figure}[H]
    \centering
    \begin{tikzpicture}[scale=0.94]
        % ---- entry points ---------------------------------------------------------------------
        \node[simple] (launcher) at (0, 0) {LauncherFrame};
        \node[simple] (simframe) at (3.6, 0) {SimulationFrame};
        \node[simple] (cli) at (7.6, 0) {Cli};
        \node[simple] (explorer) at (11.4, 0) {RuleExplorer};
        % ---- simulation -----------------------------------------------------------------------
        \node[simple] (snapshot) at (0, -1.9) {Snapshot};
        \node[simple] (simulation) at (3.6, -1.9) {Simulation};
        \node[simpleiface] (projection) at (7.6, -1.9) {\ctitle[interface]{Projection}};
        \node[simple] (tablerule) at (11.4, -1.9) {TableRule};
        % ---- config ---------------------------------------------------------------------------
        \node[simple] (library) at (0, -3.9) {DefinitionLibrary};
        \node[simple] (reader) at (3.6, -3.9) {XmlDefinitionReader};
        \node[simple] (definition) at (7.6, -3.9) {AutomatonDefinition};
        \node[simple] (factory) at (11.4, -3.9) {AutomatonFactory};
        % ---- automaton ------------------------------------------------------------------------
        \node[simple] (automaton) at (7.6, -6.0) {CellularAutomaton};
        \node[simple] (regime) at (11.4, -6.0) {Regime};
        % ---- rule -----------------------------------------------------------------------------
        \node[simpleiface] (rule) at (3.0, -6.0) {\ctitle[interface]{Rule}};
        \node[simple] (exprule) at (3.0, -7.7) {ExpressionRule};
        \node[simple] (parser) at (-0.2, -7.7) {RuleParser};
        \node[simpleiface] (expression) at (3.0, -9.4) {\ctitle[interface]{Expression}};
        \node[simple] (operator) at (0.4, -11.1) {\atitle{Operator}};
        \node[simple] (count) at (3.6, -11.1) {Count};
        % ---- neighborhood ---------------------------------------------------------------------
        \node[simple] (catalog) at (7.3, -7.7) {NeighborhoodCatalog};
        \node[simpleiface] (neighborhood) at (7.3, -9.4) {\ctitle[interface]{Neighborhood}};
        \node[simple] (offsets) at (7.3, -11.1) {OffsetNeighborhood};
        % ---- grid -----------------------------------------------------------------------------
        \node[simple] (grid) at (11.6, -7.7) {Grid};
        \node[simple] (array) at (11.6, -9.4) {\atitle{DynamicArrayND<T>}};
        \node[simpleiface] (cell) at (11.6, -11.1) {\ctitle[interface]{Cell}};
        \node[simple] (gridcell) at (10.5, -12.6) {GridCell};
        \node[simple] (voidcell) at (12.8, -12.6) {VoidCell};

        % ---- relationships --------------------------------------------------------------------
        \draw[uses] (launcher.west) -- ++(-0.45, 0) |- (library.west);
        \draw[uses] (launcher) -- node[elabel] {opens} (simframe);
        \draw[association] (simframe) -- (simulation);
        \draw[uses] (simulation) -- node[elabel] {publishes} (snapshot);
        \draw[association] (simulation) -- (projection);
        \draw[uses] (cli) -- (projection);
        \draw[uses] (cli) -- (explorer);
        \draw[association] (explorer) -- (tablerule);
        \draw[uses] (tablerule) -- node[elabel, pos=0.45] {exports} (definition.north east);
        \draw[uses] (library) -- (reader);
        \draw[uses] (reader) -- node[elabel] {creates} (definition);
        \draw[uses] (factory) -- node[elabel] {reads} (definition);
        \draw[uses] (factory) -- node[elabel, pos=0.4] {creates} (automaton.north east);
        \draw[composition] (simulation) -- (automaton.north west);
        \draw[composition] (automaton) -- (rule);
        \draw[composition] (automaton) -- (grid.north west);
        \draw[association] (automaton) -- (regime);
        \draw[implements] (exprule) -- (rule);
        \draw[composition] (exprule) -- (expression);
        \draw[uses] (parser) -- node[elabel, pos=0.55] {creates} (expression.north west);
        \draw[implements] (operator) -- (expression);
        \draw[implements] (count) -- (expression);
        \draw[composition] (operator.west) -- ++(-0.5, 0) |- node[elabel, pos=0.25] {operands} (expression.west);
        \draw[association] (count) -- (neighborhood.south west);
        \draw[aggregation] (catalog) -- (neighborhood);
        \draw[implements] (offsets) -- (neighborhood);
        \draw[uses] (neighborhood) -- node[elabel, pos=0.4] {yields} (cell.north west);
        \draw[composition] (grid) -- (array);
        \draw[composition] (array.south west) -- node[elabel, pos=0.45] {stores} (gridcell.north);
        \draw[implements] (gridcell) -- (cell);
        \draw[implements] (voidcell) -- (cell);

        % ---- packages -------------------------------------------------------------------------
        \begin{scope}[on background layer]
            \node[package, fit=(launcher)(simframe)] (p-ui) {};
            \node[package, fit=(cli)] (p-cli) {};
            \node[package, fit=(explorer)(tablerule)] (p-explore) {};
            \node[package, fit=(snapshot)(simulation)(projection)] (p-sim) {};
            \node[package, fit=(library)(reader)(definition)(factory)] (p-config) {};
            \node[package, fit=(automaton)(regime)] (p-automaton) {};
            \node[package, fit=(rule)(parser)(operator)(count)(expression)] (p-rule) {};
            \node[package, fit=(catalog)(neighborhood)(offsets)] (p-nb) {};
            \node[package, fit=(grid)(array)(cell)(gridcell)(voidcell)] (p-grid) {};
        \end{scope}
        \node[packagelabel, anchor=south west] at (p-ui.north west) {ui.swing};
        \node[packagelabel, anchor=south west] at (p-cli.north west) {cli};
        \node[packagelabel, anchor=south east] at (p-explore.north east) {explore};
        \node[packagelabel, anchor=south west] at (p-sim.north west) {simulation};
        \node[packagelabel, anchor=south west] at (p-config.north west) {config};
        \node[packagelabel, anchor=south east] at (p-automaton.north east) {automaton};
        \node[packagelabel, anchor=north west] at (p-rule.south west) {rule};
        \node[packagelabel, anchor=north west] at (p-nb.south west) {neighborhood};
        \node[packagelabel, anchor=north west] at (p-grid.south west) {grid};
    \end{tikzpicture}
    \caption{Class diagram of the project (main classes only).
        \legendarrow{extends}~inheritance,
        \legendarrow{implements}~realization,
        \legendarrow{composition}~composition,
        \legendarrow{aggregation}~aggregation,
        \legendarrow{association}~association,
        \legendarrow{uses}~dependency.}
    \label{fig:overview}
\end{figure}

Reading the diagram from the bottom:

\begin{itemize}
    \item a \code{Grid} stores its cells in a \code{DynamicArrayND}; a \code{Neighborhood} yields
          \code{Cell}s, which are either real \code{GridCell}s or the \code{VoidCell} standing for
          everything outside the grid;
    \item a rule is a tree of \code{Expression}s whose \code{Count} leaves walk neighborhoods; the
          engine only sees it through the one-method interface \code{Rule};
    \item a \code{CellularAutomaton} is composed of a grid and a rule, and reports its
          \code{Regime} (evolving, converged, cycling);
    \item an \code{AutomatonDefinition} is the content of an input file as immutable data; the
          \code{AutomatonFactory} turns it into a live automaton;
    \item a \code{Simulation} runs an automaton on its own thread and publishes \code{Snapshot}s
          to the window.
\end{itemize}

\subsection{Design patterns}

The course presents design patterns as a vocabulary for recurring solutions; \cref{tab:patterns}
lists those that the project uses, each for a reason given in the corresponding section.

\begin{table}[H]
    \centering\small
    \begin{tabularx}{\textwidth}{@{}l l L@{}}
        \toprule
        Pattern                & Where                                                                 & Purpose                                      \\
        \midrule
        Composite              & \code{DynamicArrayND}; \code{Expression}                              & Recursive structures handled
        uniformly                                                                                                                                     \\
        Interpreter            & \code{Expression.evaluate}                                            & The syntax tree of a rule evaluates itself   \\
        Template method        & \code{UnaryOperator}, \code{BinaryOperator}                           & Operand evaluation written
        once                                                                                                                                          \\
        Iterator               & \code{Neighborhood.around}                                            & A neighborhood \emph{is} a traversal         \\
        Strategy               & \code{BoundaryCondition}, \code{Projection}, \code{Rule}              & Interchangeable
        behaviors                                                                                                                                     \\
        Null object, Singleton & \code{VoidCell}                                                       & No bounds checks in rules                    \\
        Factory                & \code{AutomatonFactory}, \code{Neighborhoods}, \code{OperatorCatalog} &
        Construction separated from use                                                                                                               \\
        Builder                & \code{AutomatonDefinition.Builder}                                    & Readable construction of a nine-field record \\
        Observer               & \code{SimulationListener}                                             & The view follows the model without coupling  \\
        Command                & \code{cli.Command}; \code{Simulation}                                 & Uniform dispatch; serialized execution       \\
        \bottomrule
    \end{tabularx}
    \caption{Design patterns used in the project.}
    \label{tab:patterns}
\end{table}
